\section{The stationary realization problem}\label{sec:introduction}
The predictive geometry of an experiment is determined by tests that can actually be performed after its observed histories. Its finite-resource theory requires a second step: predictions must be propagated by an executable causal observer. A static code for a posterior distribution does not supply such an observer. In a controlled experiment the same retained label also selects the next acquisition, changing the distribution that is to be compressed. An indefinitely reused code must therefore solve a global calibration problem, not a separate quantization problem at each date.

We begin with positive experiment kernels, realizable mixtures and a finite separating family of executable future tests. The normalized predictive update is obtained from affine unnormalized report and test responses. This finite realization is an explicit hypothesis, verified from primitive kernels in the applications; finite physical dimension alone is not its proof. The resulting acquired laws depend on actual preparation, reports and actions. No Fisher tensor, covariance geometry or manifold is supplied as a foundational datum.

Our principal result constructs one stationary finite-state observer whose long-run predictive risk has the same memory order as unrestricted checkpoint compression. The full-history comparator optimizes the actions of the same prepared experiment. The proof has four parts. A common acquired mass uniformly over feedback makes the Bellman operator contract in span. The average-cost bias is strongly concave, so its Jensen defect both controls predictive loss from below and admits a second-order acquired-mass upper. A finite-dimensional correspondence on unnormalized label/preparation ensembles closes all calibration equations simultaneously. Finally, a small observer-label erasure and a separate raw physical minorization make that calibrated invariant ensemble attractive from every initial preparation.

The erasure is important in two distinct ways. It costs one retained label and an explicit risk term, and its rate controls the quantitative entrance time. It neither resets the physical experiment nor supplies a new observation. The fixed-point proof alone would establish a stationary ensemble, not convergence to it. Likewise, a fixed point with real coefficients is not a finite-bit program; the latter is established only under the additional effective certificates in Section~\ref{sec:compilation}.

For $n$ possible acquisition cuts and $M$ labels, let $J_*(n)$ be the full-history optimal risk of the protected audit protocol below. The stationary controller gives
\[
 cM^{-2/d}-\frac{C}{n}
 \ \le\ \calR_{\rm hom}(n,M)-J_*(n)
 \ \le\ C\left(M^{-2/d}+\frac{M^{2/d}}{n}\right).
\]
Consequently the bounds match for $n$ at least a fixed multiple of $M^{4/d}$, and the long-run excess is comparable to $M^{-2/d}$. The term $M^{2/d}/n$ is an entrance estimate for the constructed observer, not a proved optimal transient cost. An internally enforced exact deadline is a different architecture: it must pay for its phase information. Here an external auditor supplies the stopping request and pays for its own clock. This distinction is part of the theorem, not a convention suppressed in the notation.

Average-cost finite-state approximation, dynamic programming, and value-directed belief compression have substantial prior theories; in particular, near-optimal finite-state controllers are not new as an existence principle \cite{YuBertsekas2008,BMT2001,DKY2024,PB2002}. The assertion established here is a quantitative hard-state law through an actual globally calibrated invariant instrument ensemble, together with an all-controller lower, two primitive verifications and separately certified execution. The acquired-overlap, conditional-pooling and curvature mechanisms inherited from \cite{GTFBayes,GTFControlled} are identified in their proofs. The new stationary construction is proved independently of either realization. The theorem remains a sufficient class theorem, not a classification of arbitrary causal experiments.
